\documentclass{article}
\usepackage[margin=1in, a4paper]{geometry}
\usepackage{amsmath, amssymb, amsfonts}
\usepackage{graphicx}
\usepackage{xcolor}
\usepackage{listings}
\usepackage{booktabs}
\usepackage[utf8]{inputenc}
\usepackage[T1]{fontenc}
\usepackage{hyperref}
\hypersetup{colorlinks=true, urlcolor=blue, linkcolor=blue}
\usepackage{enumitem}
\usepackage{float}

\definecolor{codegreen}{rgb}{0,0.6,0}
\definecolor{codegray}{rgb}{0.5,0.5,0.5}
\definecolor{codepurple}{rgb}{0.58,0,0.82}
\definecolor{backcolour}{rgb}{0.99,0.99,0.99}
% Professional color palette for academic publishing
\definecolor{myblue}{cmyk}{0.8,0.4,0,0.1}
\definecolor{mygreen}{cmyk}{0.7,0,0.5,0.2}
\definecolor{myorange}{cmyk}{0,0.5,0.8,0.1}
\definecolor{mygray}{cmyk}{0,0,0,0.4}
\definecolor{arrowcolor}{cmyk}{0.9,0.5,0,0.3}
\definecolor{nodecolor}{cmyk}{0.6,0.2,0,0.05}
\definecolor{framecolor}{cmyk}{0.4,0.1,0,0.1}

\lstdefinestyle{mystyle}{
    backgroundcolor=\color{backcolour},   
    commentstyle=\color{codegreen},
    keywordstyle=\color{magenta},
    numberstyle=\tiny\color{codegray},
    stringstyle=\color{codepurple},
    basicstyle=\ttfamily\footnotesize,
    breakatwhitespace=false,         
    breaklines=true,                 
    captionpos=b,                    
    keepspaces=true,                 
    numbers=left,                    
    numbersep=5pt,                  
    showspaces=false,                
    showstringspaces=false,
    showtabs=false,                  
    tabsize=2
}
\lstset{style=mystyle}

\title{The Logistics \& Supply Chain Problem-Solving Playbook\\Chapter 70: The Minimum Spanning Tree Problem}
\author{The Logistics \& Supply Chain Playbook Series}
\date{}

\begin{document}
\maketitle

\section{The Minimum Spanning Tree Problem}

The Minimum Spanning Tree (MST) problem represents one of the most fundamental optimization challenges in network design and logistics infrastructure planning. At its core, the MST problem seeks to connect all nodes in a network using the minimum total cost while ensuring full connectivity without creating cycles. This elegant mathematical concept finds profound applications in supply chain network design, transportation route optimization, telecommunications infrastructure, and distribution center placement.

\subsection{The Scenario: Regional Distribution Network Optimization}

GlobalLogistics Inc. is expanding its operations across a new geographic region consisting of eight major cities. The company needs to establish a distribution network that connects all cities while minimizing the total infrastructure investment cost. Each potential connection between cities has an associated cost based on distance, terrain difficulty, and regulatory requirements.

The cities are: Atlanta (A), Boston (B), Chicago (C), Denver (D), El Paso (E), Fresno (F), Houston (G), and Indianapolis (H). The cost matrix represents millions of dollars required to establish direct logistics corridors between city pairs.

The challenge is to determine which connections to build such that all cities are connected through the network (either directly or indirectly) while minimizing the total investment cost. This is a classic Minimum Spanning Tree problem where cities represent vertices, potential connections represent edges, and connection costs represent edge weights.

\subsection{Tier 1: The Pen \& Paper Method (Mathematical Formulation)}

The MST problem can be formulated as an integer linear programming model that captures the essential structure of network connectivity while minimizing total cost.

\textbf{Sets and Parameters:}
\begin{align}
V &= \{1, 2, \ldots, n\} \quad \text{set of vertices (cities)} \\
E &= \{(i,j) : i, j \in V, i < j\} \quad \text{set of potential edges} \\
c_{ij} &= \text{cost of edge } (i,j) \in E
\end{align}

\textbf{Decision Variables:}
\begin{align}
x_{ij} &= \begin{cases}
1 & \text{if edge } (i,j) \text{ is selected in the MST} \\
0 & \text{otherwise}
\end{cases} \quad \forall (i,j) \in E
\end{align}

\textbf{Objective Function:}
\begin{align}
\text{Minimize} \quad Z = \sum_{(i,j) \in E} c_{ij} x_{ij}
\end{align}

\textbf{Constraints:}
\begin{align}
\sum_{(i,j) \in E} x_{ij} &= n - 1 \quad \text{(exactly } n-1 \text{ edges)} \\
\sum_{(i,j) \in E: i \in S \text{ or } j \in S, i \neq j} x_{ij} &\geq 1 \quad \forall S \subset V, S \neq \emptyset, S \neq V \quad \text{(connectivity)} \\
x_{ij} &\in \{0, 1\} \quad \forall (i,j) \in E
\end{align}

The first constraint ensures that exactly $n-1$ edges are selected, which is necessary for any spanning tree. The second constraint (subtour elimination) ensures that every proper subset of vertices has at least one edge connecting it to the remaining vertices, preventing disconnected components.

\begin{figure}[htbp]
\centering
\includegraphics[width=\textwidth]{../figures-png/line70.fig1.png}
\caption{MST Problem Graph Representation with Solution}
\end{figure}

\textbf{Concrete Example Solution:}

Consider the 8-city network with the following cost matrix (in millions):

\begin{table}[htbp]
\centering
\begin{tabular}{c|cccccccc}
& A & B & C & D & E & F & G & H \\
\hline
A & - & 12 & 8 & 6 & - & - & - & - \\
B & 12 & - & 5 & - & - & - & - & 7 \\
C & 8 & 5 & - & 9 & - & - & - & - \\
D & 6 & - & 9 & - & 4 & - & - & - \\
E & - & - & - & 4 & - & - & 3 & - \\
F & - & - & - & - & - & - & 5 & 6 \\
G & - & - & - & - & 3 & 5 & - & 8 \\
H & - & - & - & - & - & 6 & 8 & - \\
\end{tabular}
\caption{Connection Cost Matrix (Millions of Dollars)}
\end{table}

The optimal MST solution selects edges: (D,A), (B,C), (B,H), (D,E), (E,G), (G,F), (F,H) with total cost of $6 + 5 + 7 + 4 + 3 + 5 + 6 = 36$ million dollars.

\subsection{Tier 2: The Classic Heuristic (Python Implementation)}

Kruskal's algorithm provides an efficient greedy approach to finding the MST by iteratively selecting the smallest weight edge that doesn't create a cycle. The algorithm utilizes a Union-Find data structure to efficiently detect cycles.

\textbf{Algorithm Theory:}

Kruskal's algorithm works on the principle that the MST can be built by adding edges in order of increasing weight, as long as they don't create cycles. The algorithm maintains a forest of trees that gradually merge until a single spanning tree is formed.

Time Complexity: $O(E \log E)$ where $E$ is the number of edges, dominated by the sorting step.

\begin{figure}[htbp]
\centering
\includegraphics[width=\textwidth]{../figures-png/line70.fig2.png}
\caption{Kruskal's Algorithm Step-by-Step Process}
\end{figure}

\textbf{Pseudocode:}
\lstinputlisting[language=Python, caption=Kruskal's MST Algorithm]{.././codes-python/line70.py1.py}

\textbf{Complete Python Implementation:}
\lstinputlisting[language=Python, caption=Complete MST Implementation with Example]{.././codes-python/line70.py2.py}
\textbf{Concrete Example Output:}
\lstinputlisting[language=Python, caption=Algorithm Output]{.././codes-python/line70.py3.py}

\subsection{Tier 3: The Advanced Algorithm (Genetic Algorithm Implementation)}

While MST has optimal polynomial-time algorithms, genetic algorithms can be valuable for variations of the MST problem with additional constraints or for solving multiple related MST problems simultaneously. The genetic algorithm approach encodes potential solutions as chromosomes and evolves them toward optimality.

\textbf{Algorithm Theory:}

The genetic algorithm for MST uses a novel encoding scheme where each chromosome represents a set of edge selections. The algorithm maintains a population of candidate solutions and applies selection, crossover, and mutation operators to evolve better solutions over generations.

\begin{figure}[htbp]
\centering
\includegraphics[width=\textwidth]{../figures-png/line70.fig3.png}
\caption{Genetic Algorithm Process for MST Problem}
\end{figure}

\textbf{Pseudocode:}
\lstinputlisting[language=Python, caption=Genetic Algorithm for MST]{.././codes-python/line70.py4.py}

\textbf{Complete Implementation:}
\lstinputlisting[language=Python, caption=Complete Genetic Algorithm MST Implementation]{.././codes-python/line70.py5.py}

\textbf{Concrete Example Output:}
\lstinputlisting[language=Python, caption=Genetic Algorithm Output]{.././codes-python/line70.py6.py}

## Practice Questions

\textbf{Tier 1 Questions (Mathematical Formulation):}

1. **Network Design Problem:** A logistics company needs to connect 6 distribution centers with the following connection costs (in thousands): A-B: 15, A-C: 10, A-D: 20, B-C: 35, B-E: 25, C-D: 30, C-F: 15, D-F: 5, E-F: 40. Formulate this as an integer programming model and solve by hand using the cut property. What is the minimum total connection cost?

2. **Modified MST with Capacity:** Extend the basic MST formulation to include edge capacity constraints where each selected edge $(i,j)$ must satisfy $f_{ij} \leq M \cdot x_{ij}$ where $f_{ij}$ represents flow and $M$ is a large number. Write the complete mathematical model for a 5-node network.

\textbf{Tier 2 Questions (Classic Heuristics):}

3. **Prim's vs. Kruskal's Analysis:** Given a complete graph with 10 vertices where edge weights follow $w_{ij} = |i-j|^2$, implement both Prim's and Kruskal's algorithms. Compare their execution steps and verify they produce the same MST. What is the total cost difference between the MST and the complete graph?

4. **Dynamic MST Updates:** Starting with the MST from the 8-city example (total cost 36), suppose the cost of edge B-H increases from 7 to 15. Implement an efficient algorithm to update the MST. Does the MST structure change, and what is the new total cost?

\textbf{Tier 3 Questions (Advanced Algorithms):}

5. **Multi-Objective MST:** Design a genetic algorithm that simultaneously minimizes total cost and maximizes network reliability (defined as the minimum edge weight in the tree). Use a weighted fitness function $f(x) = w_1 \cdot \frac{1}{\text{cost}} + w_2 \cdot \text{min\_edge\_weight}$. Test with weights $w_1 = 0.7, w_2 = 0.3$ on a 12-node random graph.

6. **Constrained MST with Forbidden Edges:** Modify the genetic algorithm to handle scenarios where certain edges cannot be selected due to regulatory restrictions. If edges A-D and G-F are forbidden in the 8-city example, what is the new optimal MST cost and structure? Compare the genetic algorithm's performance against a constraint-aware Kruskal's variant.

\end{document}
